\documentclass[11pt]{article}
\usepackage[margin=1in]{geometry}
\usepackage{amsmath,amssymb}
\usepackage{xurl}
\title{Discrete-context transposition: claim-to-Lean crosswalk}
\date{20 November 2026}
\begin{document}
\maketitle

\paragraph{Audit and boundary.}
This formalization covers the discrete-transposition obstruction paragraph of
issue \#105 and the negative theorem in \emph{Wreath-extension candidates},
\S 8, at source revision
\nolinkurl{458ff36a4798e432dee5947afc402f9903a47f14}.
It proves a necessary condition on the stated hom-level interface, then
exhibits a particular monad that violates it. It is not an existence or
reconstruction theorem for contextads, a statement about arbitrary
Kleisli-category equivalence with no pure-map condition, an obstruction to
non-discrete or quotient presentations, a categorical-size generalization,
or a proof of the positive Appendix A.4 result.

\paragraph{Set/Type convention and actual maps.}
The categorical base is mathlib's \texttt{Type u}. Small sets are represented
by types in that universe and their morphisms by actual functions with the
\texttt{TypeCat} wrapper. Tags and domains are in the same universe as the
base. There are no finiteness or inhabitedness premises for that generic
claim; the finite witness is on \texttt{Type 0}. For each fixed source
$A$, the presentation has an equivalence
\[
 \mathrm{Hom}(A,T B)\simeq \sum_{p:E(A)}\mathrm{Hom}(D(A,p),B)
\]
at every $B$, and one law for every actual pure arrow induced by $f:B\to C$.
On the left this is composition by $\eta_C f$ in the Kleisli category; on
the right it is $(p,h)\mapsto(p,f h)$, with no quotient or change of tag.
No other contextad operation is imposed to derive the obstruction.

\section*{Generic chain}
The following names are under
\nolinkurl{ContextualComposition.DiscreteTransposition}.

\begin{description}
\item[Both postcomposition actions.]
\nolinkurl{Presentation}, \nolinkurl{polynomial_map},
\nolinkurl{purePost_eq_map}, and \nolinkurl{Presentation.homEquiv_natural}.
\nolinkurl{ContextualComposition.MonadRecovery.kleisli_comp} is the inherited
native-composition identification. Mathlib's monad right-unit law reduces
$\mu_C\,T(\eta_C f)\,k$ to $T(f)\,k$.
\item[Singleton extraction.]
\nolinkurl{singletonEval} and \nolinkurl{Presentation.singletonIso}.
Constant sections and evaluation give mutually inverse maps at $A=1$.
The component square uses \nolinkurl{Presentation.homEquiv_natural}; the
result is a bundled natural isomorphism to \nolinkurl{polynomial}, not just
objectwise equality of cardinalities or equivalence of carriers.
\item[Arbitrary cospan.]
\nolinkurl{MappedPullback}, \nolinkurl{imageCone},
\nolinkurl{comparison} and \nolinkurl{polynomial_comparison_bijective}.
Both projections come from mathlib's explicit limit pullback cone for
\texttt{Type}. The target is the subtype of pairs agreeing under the
mapped legs. Equality in the coproduct forces equality of tags. With equal
tags the lifts pair pointwise over their common domain, and equality on both
projections implies equality of the original lifts.
\item[Transport and categorical reading.]
\nolinkurl{comparison_bijective_of_iso},
\nolinkurl{Presentation.comparison_bijective}, and
\nolinkurl{Presentation.imageConeIsLimit}.
Naturality on both legs ensures that the transported pair is in the target.
Naturality on both projections makes the comparison square commute.
The explicit image cone is limiting by mathlib's equivalence between that
property and bijectivity of the canonical comparison.
\end{description}

\section*{Continuation witness}
The following names are under
\nolinkurl{ContextualComposition.ContinuationTranspositionObstruction}.

\begin{description}
\item[Actual monad and square.]
\nolinkurl{continuationFunctor}, \nolinkurl{continuationMonad},
\nolinkurl{continuation_map}, \nolinkurl{toUnit}, and
\nolinkurl{terminal_map}. The action is $k\mapsto(q\mapsto k(qf))$;
unit is evaluation and multiplication is
$K\mapsto(q\mapsto K(k\mapsto k(q)))$. Functoriality, the two
naturality equations and the three monad equations are proved. Both legs
are $\mathrm{Bool}\to\mathrm{Unit}$.
\item[Exact source and target.]
\nolinkurl{source_card}, \nolinkurl{terminal_card},
\nolinkurl{fiberEquiv}, \nolinkurl{fiber_card},
\nolinkurl{targetFibers}, and \nolinkurl{target_card}.
The source is the continuation type of the four-point pullback, so it
has $2^{2^4}=65536$ elements. A continuation on the singleton fixes exactly
the values on the two constant Boolean predicates. Its lifts freely set
the values on the identity and negation predicates, exhibited by an
explicit inverse for each fiber. The target subtype splits over the four
common images, giving $4\cdot4\cdot4=64$.
\item[Obstructions.]
\nolinkurl{comparison_not_bijective}, \nolinkurl{no_polynomial},
and \nolinkurl{no_presentation}. A bijection of the actual comparison
would identify the two cardinalities. This rules out a natural
coproduct-of-representables presentation, hence the given discrete
transposition. It does not use cardinalities of an abstract substitute
for either endpoint of the comparison.
\end{description}

\paragraph{Discriminating checks.}
\nolinkurl{checks/DiscreteTransposition.lean} verifies an invertible Boolean
chart for the identity monad whose square fails under a particular
nonidentity pure map; the chart example is local and does not deny that
the identity monad has a different lawful presentation. It also checks
the two different actual projection outputs on a pure evaluation, rejects
an ambient pair whose constant-predicate values disagree, and witnesses
an actual collision on the source comparison using a mixed predicate
invisible to either projection separately. These checks do not rely on
the cardinal arithmetic and do not replace the generic proof.

\paragraph{Provenance and replay.}
The inherited projection is locked to 507 source paths of merged parent
\nolinkurl{f458186590df968a516ba73a8eb76fcf044544b0},
tree \nolinkurl{1c829c869d1c15e907b1f2404cad74e6fd88f1b7}.
The target toolchain is Lean 4.35.0 with mathlib
\nolinkurl{56a005855ea17e0b0958f73ec801ae140b0a1679} and all transitive
pins in the inherited Lake manifest. The change adds two production
modules and one probe, and updates the umbrella and paired audit/index.
Warm author replay imports copied, digest-verified inherited products.
A cold replay must start with no project products and rerun the complete
ordered audit and all probes. The repository evidence includes a cold replay runner
and exact input manifests. The warm run is not the independent replay.
\end{document}
